% =====================================================================
% Section 5 — Discussion
% =====================================================================

\section{Discussion}\label{sec:discussion}

\subsection{Implications for cylinder-head design}
\label{sec:disc-design}

The quantitative-coupled BowTie of Section~\ref{sec:results} provides a
single artefact that links each design decision in the high-performance
cylinder head to the risk it controls. The analysis identified two
concurrent at-risk barriers under the baseline configuration of the
AP 1.8 head: (i) the \emph{outer-spring coil-bind margin}
($M_{\mathrm{out}}=0.78$~mm against a project criterion of
$M\geq 1.5$~mm), and (ii) the \emph{inner-spring fatigue Goodman
factor of safety} ($n_{f,\mathrm{in}}\approx 0.68$ against
$n_{f}\geq 1$). The outer-spring natural-frequency margin
($f_{n,\mathrm{out}}/f_{\mathrm{cam,max}}\approx 3.81$) sits just
below the conventional rule-of-thumb threshold of 4 and is treated
as a residual concern. The remaining barriers --- dynamic follower
contact ($R\geq 1.343$ in the worst case, against the project
criterion of $R\geq 1.20$), inner-spring coil-bind margin
($M_{\mathrm{in}}=2.38$~mm), and outer-spring fatigue
($n_{f,\mathrm{out}}\approx 1.27$) --- satisfy their pass criteria
with non-uniform margins.

The fact that the two at-risk barriers affect \emph{different}
springs (outer for coil-bind, inner for fatigue) and act through
\emph{different} physical mechanisms (geometric clearance vs cyclic
shear loading) is significant: a design iteration that addresses
only one of the two findings will leave the other unresolved. The
sizing note for the AP 1.8 springs already prescribes the geometric
correction for the outer-spring coil-bind (raising the installed
height to $H_{\mathrm{inst,ideal}}=47.72$~mm and the free length to
$L_{0,\mathrm{ideal}}=55.56$~mm to achieve $M\approx 2.0$~mm). The
inner-spring Goodman deficit, by contrast, calls for a metallurgical
or geometric change of a different nature: either increasing the
wire diameter $d_{\mathrm{in}}$ to lower the shear stress at a given
load, reducing the spring index
$C_{\mathrm{in}}=D_{m,\mathrm{in}}/d_{\mathrm{in}}$, or specifying
a higher-grade steel (e.g.\ chrome--silicon--vanadium ASTM~A878)
with a larger corrected shear endurance limit $S_{se}$. Concerning
the residual outer-spring natural-frequency concern, a parallel
correction would adjust either the active coil count
$N_{a,\mathrm{out}}$ or the spring index $C_{\mathrm{out}}$ to push
$f_{n,\mathrm{out}}$ further from the dominant cam harmonics. The
BowTie diagram makes the three trade-offs simultaneously visible to
the design team.

The remaining barriers retain non-zero margins and therefore satisfy
the literature-anchored pass criteria, but Table~\ref{tbl:quant-results}
shows that the margins are not uniform. This non-uniformity is a useful
input for design iteration: barriers with comfortable margins (e.g.\
the inner-spring coil-bind margin or the outer-spring fatigue stress)
are candidates for relaxation if doing so frees mass or simplifies
manufacturing; barriers close to the pass threshold should be the focus
of dynamometer validation as soon as a physical prototype is available.

% ----------------------------------------------------------------------
\subsection{Comparison with traditional FMEA/RPN}
\label{sec:disc-fmea-comparison}

Table~\ref{tbl:bowtie-vs-fmea} contrasts the present BowTie analysis
with a conventional FMEA of the same six threats. Three observations
follow.

\begin{enumerate}
\item \emph{Granularity of the causal chain}: the BowTie diagram makes
the linkage between threat, preventive barrier, top event, and
consequence visible in a single artefact, whereas FMEA collapses the
causal chain into the multiplicative RPN. For a high-rpm valve train,
this granularity is necessary because several threats share the same
preventive barriers (for example, both \emph{insufficient stiffness}
and \emph{spring heating} are addressed by the spring-natural-frequency
barrier).
\item \emph{Quantitative anchoring}: in FMEA, severity, occurrence, and
detectability are typically assigned on integer scales~(1--10) by
expert judgement. The quantitative-coupled BowTie replaces these
ordinal scores with deterministic indicators that can be reproduced by
any reader who has access to the sizing calculations.
\item \emph{Communication}: the BowTie diagram is accessible to
designers, operators, and reliability engineers without prior
familiarity with the RPN convention. This is a practical advantage in
the academic context of a master's research that aims to articulate
mechanical design and reliability analysis.
\end{enumerate}

\begin{table}[t]
\caption{Comparison between the BowTie analysis presented in this work and a conventional FMEA of the same threats.}\label{tbl:bowtie-vs-fmea}
\footnotesize
\begin{tabular*}{\tblwidth}{@{}LCC@{}}
\toprule
Aspect & FMEA / RPN & Quantitative BowTie \\
\midrule
Causal chain & Implicit & Explicit \\
Severity assessment & Ordinal (1--10) & Continuous (deterministic) \\
Pass criterion & RPN threshold & Indicator-by-indicator \\
Reproducibility & Limited & Direct from sizing \\
Visualisation & Tabular & Diagrammatic \\
Identifies critical barrier & No & Yes (via tornado plot) \\
\bottomrule
\end{tabular*}
\end{table}

A concrete numerical illustration of why the conventional RPN
collapses information that the BowTie preserves is given in
Table~\ref{tbl:rpn-comparison}. The six threats $T_{1}$--$T_{6}$ of
the BowTie diagram were scored on the conventional 1--10 ordinal
scales of severity ($S$), occurrence ($O$), and detectability ($D$)
following the rating guidelines of the AIAG--VDA 2019 FMEA harmonised
handbook \citep{aiagvda2019fmea}: severity is uniformly high
because all six threats converge on the same catastrophic top event
(valve--piston collision, $S\in[8,10]$); occurrence is mapped from
the ASTM~A877/A877M tolerance band and the catalogued operational
variability of the spring (so that occurrence is calibrated to
manufacturing rather than to subjective judgement); detectability is
uniformly low ($D\in[3,4]$) because the failure modes are mechanical
and predominantly silent before the event of failure. Each score is
the median of three independent assignments by the authors using the
AIAG--VDA rubric. The resulting RPNs fall in a narrow band
(96--240) and place the inner-spring fatigue threat ($T_{4}$,
RPN~=~150) and the outer-spring coil-bind threat ($T_{1}$,
RPN~=~144) within ten points of each other, even though the
quantitative-coupled BowTie identifies both as concurrently at risk
through entirely different physical mechanisms. A triaging team
using a conventional 100--125 RPN threshold would not be able to
distinguish reliably between these two threats. The RPN ranking, in
other words, masks the very finding that the BowTie analysis
surfaces. Recent extensions of FMEA --- the action-priority logic
introduced in the AIAG--VDA 2019 handbook itself, fuzzy-set
formulations \citep{liu2013fuzzyfmea,wang2009fuzzyAHP}, and
AHP-weighted scoring --- improve the resolution of the prioritisation
but do not address the underlying compression of the causal chain
into a single ordinal score, which is the structural objection
addressed by the present BowTie formulation.

\begin{table}[t]
\caption{AIAG--VDA 2019 scoring of the six threats $T_{1}$--$T_{6}$ of the BowTie diagram. Severity, occurrence, and detectability follow the harmonised severity-occurrence-detection rubric of \citet{aiagvda2019fmea}; each score is the median of three independent assignments by the authors. Occurrence is calibrated to the ASTM~A877/A877M tolerance band rather than to subjective frequency judgement.}\label{tbl:rpn-comparison}
\begin{tabular*}{\tblwidth}{@{}LCCCC@{}}
\toprule
Threat & $S$ & $O$ & $D$ & RPN \\
\midrule
$T_{1}$ Insufficient spring stiffness & 9 & 4 & 4 & 144 \\
$T_{2}$ Excessive moving mass & 9 & 3 & 4 & 108 \\
$T_{3}$ Aggressive cam profile & 10 & 3 & 4 & 120 \\
$T_{4}$ Spring heating in operation & 10 & 5 & 3 & 150 \\
$T_{5}$ Out-of-tolerance thermal clearance & 8 & 4 & 3 & 96 \\
$T_{6}$ Cam-harmonic / spring resonance & 10 & 6 & 4 & 240 \\
\bottomrule
\end{tabular*}
\end{table}

The two methods are not exclusive: FMEA remains useful for triaging a
large number of failure modes when probabilistic data is unavailable,
and BowTie is most valuable when a small number of high-severity events
deserves a detailed barrier analysis. For the high-performance cylinder
head studied here, the BowTie approach is preferred because the failure
modes converge on a single top event whose consequences are catastrophic.

A parallel branch of the literature has extended the conventional FMEA
through fuzzy-set formulations and AHP-based weighting to address the
loss of resolution of the multiplicative RPN. Compared with these
approaches, the proposed quantitative-coupled BowTie does not refine
the FMEA score; it \emph{replaces} the single-score logic with a
barrier-by-barrier deterministic check that is reproducible from the
sizing calculations. The two routes are therefore complementary:
fuzzy/AHP-FMEA improves the prioritisation among many failure modes,
while the present BowTie improves the auditability of the design
margins around a single critical event.

% ----------------------------------------------------------------------
\subsection{Limitations}\label{sec:disc-limitations}

The present analysis has the following limitations, which should be
considered when interpreting the results. They are listed in order
of priority for follow-on work.

\begin{itemize}
\item \emph{No experimental validation --- primary limitation.} The
quantitative indicators of Table~\ref{tbl:quant-results} are
computed from the analytical sizing of the valve and the spring;
no dynamometer, motored-engine, or spring-rig test of the
high-performance cylinder head has been conducted. The qualitative
validation against four published failure cases reported in
Section~\ref{sec:res-litvalidation} supports the diagnostic
behaviour of the procedure but does \emph{not} confirm the absolute
magnitude of the indicator margins for the AP 1.8 head. An
empirical campaign is required, and its scope is now stated
quantitatively for transparency. The minimum experimental
verification programme that would close this limitation comprises:
(i) a dedicated spring rig to measure the surge response of both
springs across the 50--600~Hz band, with laser-displacement
instrumentation on the active coils, sufficient to confirm
$f_{n,\mathrm{out}}\approx 318$~Hz and
$f_{n,\mathrm{in}}\approx 571$~Hz to within $\pm$5\,\% under the
seat-preload boundary condition; (ii) a motored-engine speed sweep
from 6\,000 to 10\,000~rpm with a non-contacting valve-position
transducer sampled at 50\,kHz, to verify the absence of valve
float at the deceleration peak ($a_{\max,-}=20\,362$~m\,s$^{-2}$)
and to measure the residual coil-bind margin within $\pm$0.1~mm;
and (iii) a fatigue campaign on a representative inner-spring
specimen at the alternating shear stress of
$\tau_{a,\mathrm{in}}\approx 374$~MPa superimposed on the mean
$\tau_{m,\mathrm{in}}\approx 858$~MPa, run to either run-out at
$10^{8}$ cycles or first-failure indication, on a sample size of
at least $n=5$ to support a statistically meaningful estimate of
the Goodman factor of safety against the
$n_{f,\mathrm{in}}\approx 0.68$ analytical prediction. The total
test-rig hours required for this minimum programme are estimated
at the order of 600--800~h. The validation roadmap is presented in
detail in Section~\ref{sec:conclusion}.
\item \emph{Operating envelope partially characterised.} The
deterministic pass criteria are evaluated at the maximum design
speed of 10\,000~rpm, and Section~\ref{sec:res-envelope} extends
the two speed-sensitive indicators ($R$ and
$f_{n}/f_{\mathrm{cam}}$) across 6\,000--12\,000~rpm at fixed
geometry. Off-design behaviour involving thermal transients
(cold start, sustained high-coolant-temperature operation) and
load transients (part-load, overspeed beyond 12\,000~rpm) is not
considered, and a coupled thermal--mechanical sweep is identified
as a follow-on analytical study.
\item \emph{Probabilistic refinement covers manufacturing tolerance
but not full operational variability.} The deterministic indicators
have been complemented by a Monte Carlo simulation
(Section~\ref{sec:res-mc}) propagating the catalogued ASTM~A877/A877M
spring-wire tolerances and the strength dispersion of
\citet[\S~10-30]{shigley2020} through the analytical sensitivity
coefficients. The simulation confirms that the inner-spring Goodman
deficit is present in 100\,\% of the manufactured population, the
outer-spring coil-bind margin fails its criterion in 77.9\,\% of
trials with a 5\,\% probability of actual hard bind, and the
inner-spring coil-bind margin is re-classified from active
(deterministic) to marginal ($P[\mathrm{fail}]=16.2\,\%$). However,
the Monte Carlo treats only the manufacturing tolerance and the
strength scatter; operational variability beyond ambient temperature
(thermal soak transients, oil-condition degradation, lubricant
contamination, valve-seat wear) and the cam-profile harmonics
distribution are not propagated. A second-generation analysis that
extends the simulation to those operational sources and computes
joint barrier-failure probabilities (e.g.\ correlated coil-bind
and Goodman events under thermally degraded conditions) is identified
as a follow-on direction.
\item \emph{Hot clearance treated as a 1-D differential estimate.}
The hot-clearance indicator
$c_{\mathrm{cold,req}}=1.25$~mm (bracket [1.13, 1.60]) of
Section~\ref{sec:res-quant} is a closed-form differential of the
steel-stem and aluminium-head expansion paths along the
load-bearing line that closes the rocker--valve-tip clearance. It
relies on a documented stem temperature
$\Delta T_{\mathrm{stem}}\approx 600$~K
\citep{heisler2002}, a head temperature
$\Delta T_{\mathrm{head}}\approx 150$~K
\citep{heywood1988}, and a generic OHC head path length
$L_{\mathrm{Al}}\approx 100$~mm. A coupled
thermal--structural finite-element analysis of the specific
AP 1.8 cylinder head would refine $L_{\mathrm{Al}}$ and the
localised $\Delta T_{\mathrm{head}}$ field and tighten the
[1.13, 1.60]~mm bracket, but the present deterministic indicator
is no longer a qualitative verification.
\item \emph{Barrier set restricted to the spring--cam dynamic
interaction.} The qualitative validation of
Section~\ref{sec:res-litvalidation} flagged two failure mechanisms
that are not currently covered by the preventive set:
out-of-tolerance valve-seat alignment
\citep{soffritti2018wornvalvetrain} and lubrication-regime
degradation in the camshaft bearing pair
\citep{zhao2023camshaftV6}. Extensions of the procedure with
candidate barriers on assembly tolerance and on the lubrication
regime are identified as priority directions for future work to
broaden the coverage of valve-train-related failure mechanisms.
\item \emph{Single design instance.} The analysis is conducted on a
single bespoke cylinder head ($n=1$). The methodology itself is
transposable to other high-rpm SI heads, and
Section~\ref{sec:res-envelope} demonstrates the procedure across the
operating envelope 6\,000--12\,000~rpm for the same head, recovering
in closed form the surge, dynamic, and loss-of-contact boundaries.
The numerical findings (critical barrier identity, indicator
margins) at the 10\,000~rpm design point are nevertheless
case-specific and should not be generalised to other heads
without re-evaluation of the geometric and material inputs.
\item \emph{Out of scope by design choice.} Knock and other abnormal
combustion phenomena, low-speed pre-ignition, IoT-based monitoring,
artificial-intelligence-supported diagnostics, and the multi-objective
optimisation of the cylinder-head geometry are explicitly out of scope.
The reader is referred to the underlying dissertation
\citep{porto2026dissertation} for the air-flow and combustion analysis
of the same head.
\end{itemize}
